\begin{frame}{손으로 계획 한 스텝: 알려진 모델에 가치반복}
  모델이 \textbf{완벽히 알려져} 있을 때 ``계획''은 어떤 모습인가 ---
  가장 단순한 형태인 \kb{가치반복}(Week 4.3)을 손으로 들여다본다.
  \vskip0.5em
  \begin{itemize}
    \item 3상태 예제: $S$(출발) $\to$ $X$(중간) $\to$ $T$(목표, 종료),
      각 상태에서 행동 하나씩, 전이는 결정론적
  \end{itemize}
  \[
  S \xrightarrow{a} X \;(\text{보상 } 0), \qquad
  X \xrightarrow{a} T \;(\text{보상 } +1), \qquad \gamma = 0.9
  \]
  \vskip0.5em
  \begin{columns}[T]
    \begin{column}{0.48\textwidth}
      \kb{참값 (거꾸로 대입)}
      \\[0.4em]
      {\footnotesize $Q(T)=0$,
      $Q(X)=1+0.9 \cdot 0 = 1$,
      $Q(S)=0+0.9 \cdot 1 = 0.9$ --- 손에 바로 나온다.}
    \end{column}
    \begin{column}{0.48\textwidth}
      \kb{가치반복 한 반복}
      \\[0.4em]
      {\footnotesize 모든 $(s,a)$에 대해
      $Q(s,a) \leftarrow r + \gamma \max_{a'} Q(s',a')$.
      주의: \textbf{동기화(synchronous) 갱신} --- 새 값을
      \textbf{오래된 Q}로 계산한 뒤 한꺼번에 바꿔 넣는다.}
    \end{column}
  \end{columns}
\end{frame}
